\section{Introduction}

The DW-NOMINATE data we downloaded shows a clear polarization trend:
mean $|\text{DW-NOMINATE dim1}|$ for House Republicans was 0.428 in the 2000s,
0.483 in the 2010s, and 0.505 in the 2020s.
Democrats showed a smaller but real increase: 0.371, 0.386, and 0.382.

Three mechanisms link redistricting to this polarization:

\begin{enumerate}
  \item \textbf{Safe-seat creation}: partisan gerrymandering creates districts
    where the primary is the effective election.
    Primary electorates are more ideologically extreme than general electorates
    \citep{brady2007move}, producing more extreme winners.
  \item \textbf{Geographic sorting}: Democrats have become more concentrated
    in urban areas, creating naturally safe Democratic urban districts
    regardless of gerrymandering \citep{rodden2019cities}.
    Algorithmic redistricting cannot undo geographic sorting but can avoid
    amplifying it via cracking and packing.
  \item \textbf{Incumbent sorting}: incumbents from safe seats, knowing they
    face no general-election threat, take more extreme positions to maintain
    primary standing \citep{canes-wrone2002out}.
\end{enumerate}

This paper focuses on mechanisms (1) and (2): the geometric relationship between
district compactness and member ideology.
If gerrymandering-induced safe seats drive extremism, then more compact districts
--- which are less likely to be deliberately packed --- should produce less extreme
members, conditional on the district's baseline partisan composition.
